\documentclass[conference]{IEEEtran}
\usepackage[T1]{fontenc}
\usepackage{url}
\begin{document}
\title{Evaluating Deep Reinforcement Learning for Traffic Signal Control Under Real Toronto Demand}
\author{\IEEEauthorblockN{Jowel Akkeh}\IEEEauthorblockA{Department of Civil Engineering, Lassonde School of Engineering\\York University, Toronto, Canada}}
\maketitle
\begin{abstract}This check-in reports an implemented multimodal SUMO baseline, a frozen train/dev/test split, and a bounded DQN/PPO pilot. Fixed-time dev vehicle delay proxy is 166.21 $\pm$ 10.16 s/vehicle. DQN reproduces the baseline; PPO is unstable and worse on average. No robust DRL benefit is established. Hosted evidence links remain a submission dependency.\end{abstract}
\section{METHOD}
\paragraph{Data and split} The implemented loader selects TCS 379, Eglinton Avenue East / Danforth Road, from Toronto's turning-movement counts [1]. Cars, trucks, and buses retain their recorded movements; pedestrian fields identify the leg being crossed. Counts are validated and reproduced by 15-minute interval. Before controller experiments, record identifiers and file hashes were saved. Table I gives the frozen partition. The original June afternoon training period was divided into a two-hour train interval and a one-hour dev interval. Both proposed test periods remain reserved: only demand totals were audited, not controller performance. No test routes or scores were generated by the learning pipeline.
\begin{table}[ht]\centering\footnotesize
\caption{FROZEN OBSERVED DEMAND PARTITIONS (2024)}
\begin{tabular}{lrrr}\hline
Set / date / period & Rows & Veh. & Peds \\
Train: Jun. 19, 15-17 & 8 & 6,716 & 810 \\
Dev: Jun. 19, 17-18 & 4 & 3,594 & 518 \\
Test: Nov. 2, 15-18 & 12 & 9,961 & 1,598 \\
Test: Jun. 19, 07-10 & 12 & 8,275 & 661 \\
\hline\end{tabular}\end{table}
\paragraph{Simulator and baseline} SUMO 1.28.0 uses a georeferenced OpenStreetMap road extract [2], with lane assignments and movement connections adapted for this isolated junction. It includes sidewalks, four full-width signalized crossings, and separate eastern carriageways. Lane widths, corner radii, sidewalk dimensions, the 2.4 m eastern median, transit operations, and historical 2024 fidelity remain uncalibrated. The fixed-time baseline uses a simplified interpretation of the published TCS 379 PM card [3]: advance green/yellow/red of 7/3/1 s, EW WALK/clearance of 31/22 s, yellow/red of 4/3 s, NS WALK/clearance of 8/34 s, then yellow/red of 4/3 s. Total cycle length is 120 s; detector calls are not reproduced.
\paragraph{Demand realization} Exactly the observed vehicle and pedestrian totals are generated. Arrivals are uniform within each observed interval; pedestrian directions are approximately balanced because directional counts are absent. Walk speed is fixed at 1.2 m/s, with trip endpoints 0.5 m from the junction. These are explicit assumptions, not measured trajectories. Each episode starts empty, has no warm-up, and permits up to 3,600 s of drainage after arrivals end. Controllers share the same geometry and seeded arrival realizations. Pending vehicle entry is retained in the delay accounting.
\paragraph{Learning environment} Stable-Baselines3 DQN and PPO [4] use a custom Gymnasium adapter with the TraCI-compatible libsumo API [5]. The installed SUMO-RL default phase builder does not retain this complete clearance sequence, so its phase-switching controller is not used. One action per cycle selects EW WALK from \{7, 19, 25, 31, 32\} s; NS WALK receives the remainder of 39 s. Both WALK periods remain at least 7 s. All other intervals and the phase order are immutable. The observation contains per-leg vehicle counts, queues, pending insertions, waiting pedestrians, maximum pedestrian wait, and the previous action. Scaling is fixed; no statistics are fitted on dev.
\paragraph{Reward and training} The reward is the negative integral of approach vehicle queues, vehicles awaiting insertion, and stopped pedestrians, with equal road-user weights. It is sampled every 5 s and divided by 1,200 per cycle; unfinished demand receives an additional terminal penalty. Each algorithm was trained for 2,048 cycle decisions for seeds 101, 202, and 303, using only train demand. Each used a 64-by-64 MLP, learning rate 0.0003, discount 0.99, and batch size 64. DQN used a 20,000-entry replay buffer, 180-decision learning start, and final exploration 0.05. PPO used 256-step rollouts and 10 epochs. Full settings were recorded before learning. One preset configuration per algorithm and one documented fixed plan were evaluated; no hyperparameter search was performed.
\paragraph{Metrics and hardware} Vehicle timeLoss measures loss relative to ideal travel speed; departDelay measures waiting before insertion [6]. Their sum is reported as a vehicle delay proxy, alongside the two components. Pedestrian walk timeLoss is reported separately. These simulation proxies are not field-calibrated intersection delay. Completion, collisions, teleports, and emergency braking are retained; incomplete demand would invalidate an all-demand delay ranking. Runs used Windows, Python 3.14.6, an Intel Family 6 Model 154 CPU with 20 logical processors, PyTorch 2.15.0+cpu, and libsumo 1.28.0. No GPU was used.
\section{FIRST EXPERIMENTS}
\paragraph{Baseline and checks} The data loader and fixed-time simulation run end to end. Across five dev demand seeds (11, 22, 33, 44, 55), mean vehicle time loss was 55.12 $\pm$ 3.72 s/vehicle and departure delay was 111.09 $\pm$ 9.22 s/vehicle. Combined delay was 166.21 $\pm$ 10.16 s/vehicle, while pedestrian time loss was 46.93 $\pm$ 1.43 s/person. Values are mean $\pm$ sample SD. All 3,594 vehicles and 518 pedestrians finished in each run. The new adapter reproduced the saved baseline exactly. A 600-state action audit checked every legal allocation, and a 7,200-state trace checked 60 fixed cycles and all intended pedestrian crossings.
\paragraph{Pilot comparison} Each last checkpoint was evaluated deterministically on the same five dev realizations: 15 cases per RL algorithm and five fixed-time cases. Table II gives case means; Table III reports policy-seed variation. DQN selected the baseline 31/8 allocation for every dev cycle across all three trained policies and therefore reproduced its results exactly. PPO was unstable across training seeds: its vehicle-proxy means were 253.89, 227.53, and 152.52 s/vehicle. Its average was worse than fixed-time. The favorable seed is not selected as proof of improvement. This initial pilot does not establish a DRL benefit.
\begin{table}[ht]\centering\footnotesize
\caption{DEV CASE MEANS (SECONDS PER ROAD USER)}
\begin{tabular}{lrrrr}\hline
Control & Veh. loss & Entry & Combined & Ped. loss \\
fixed & 55.12 & 111.09 & 166.21 & 46.93 \\
DQN & 55.12 & 111.09 & 166.21 & 46.93 \\
PPO & 58.76 & 152.55 & 211.31 & 46.61 \\
\hline\end{tabular}\end{table}
\paragraph{Uncertainty and limits} A paired, two-way bootstrap resampled three policy seeds and five shared demand seeds (4,000 samples; random seed 5420). PPO minus fixed-time combined vehicle delay was +45.11 s, with exploratory 95\% interval [-12.28, 87.64] s. The pedestrian difference was -0.32 s, interval [-0.92, 0.24] s. Both span zero; small improvements are not established. DQN differences are identically zero because its actions match the baseline. No collisions or teleports occurred in the 35 dev cases; emergency braking averaged 1.2, 1.2, and 2.0 events per case for fixed-time, DQN, and PPO. Repeated stochastic episodes do not add independent observed days. Time-blocked cross-validation has not yet been performed.
\begin{table}[ht]\centering\footnotesize
\caption{EACH POLICY: FIVE DEMAND SEEDS, MEAN $\pm$ SD}
\begin{tabular}{lrr}\hline
Policy seed & Vehicle combined & Pedestrian loss \\
PPO 101 & 253.89 $\pm$ 14.82 & 46.24 $\pm$ 1.65 \\
PPO 202 & 227.53 $\pm$ 12.92 & 46.54 $\pm$ 1.54 \\
PPO 303 & 152.52 $\pm$ 7.62 & 47.07 $\pm$ 1.41 \\
DQN: all 3 & 166.21 $\pm$ 10.16 & 46.93 $\pm$ 1.43 \\
\hline\end{tabular}\end{table}
\section{EVIDENCE, ISSUES, AND REVISED PLAN}
\paragraph{Reproducibility and hosted evidence} The local repository contains the loader, frozen row identifiers and hashes, model, control adapter, six trained checkpoints, per-seed metrics, and configuration. The experiment is under experiments/pilot\_v1; its dev\_summary.json is the numerical source for this report. W\&B recorded six training runs with learning histories and model artifacts, plus a matched-dev comparison. These are offline runs indexed by wandb\_runs.json. Hosted repository URL: not yet created. Hosted W\&B URLs: not yet synchronized. Account sign-in and publication remain required to satisfy the evidence-link checklist; local artifacts are not presented as hosted links.
\paragraph{Problems found and corrected} Runtime tracing exposed a guessed crossing that let 62 south-leg pedestrians bypass the signal; that shortcut was removed and all 518 dev pedestrians were verified on the correct legs. A preliminary training attempt failed because a diagnostic and training process shared a temporary statistics file. That attempt was archived and excluded; all successful models restarted after process-specific output isolation. These issues and the null/unstable learning result are retained as evidence rather than hidden.
\paragraph{Revised scope} The check-in now has a running baseline and a bounded pilot; further training is not required to make this report concrete. Next, publish the repository and synchronize the existing W\&B runs, then insert verified links. A later development experiment should inspect PPO's green allocations and entry queues, consider a separately tuned fixed-allocation diagnostic with comparable dev tuning effort, and revise training budget/reward only through dev. Field calibration and time-blocked validation remain necessary before broad claims. Test periods stay sealed until the final configuration is frozen; changing a model after viewing test performance would require relabeling that test as dev.
\begin{thebibliography}{6}
\bibitem{r1} City of Toronto, "Traffic volumes - multimodal intersection turning movement counts," Toronto Open Data. Dataset and dictionary, accessed Nov. 11, 2026.
\bibitem{r2} OpenStreetMap contributors, "Eglinton Avenue East / Danforth Road map extract," ODbL 1.0, downloaded Nov. 9, 2026. Attribution and licence.
\bibitem{r3} City of Toronto, "Current Signal Timing Information: TCS 379," Oct. 31, 2019, in WSP, 320 McCowan Road Transportation Impact Study, Feb. 14, 2020, App. B, PDF p. 147. Source report.
\bibitem{r4} Stable-Baselines3, "Using custom environments," v. 2.9.2 documentation. Documentation, accessed Nov. 11, 2026.
\bibitem{r5} Eclipse SUMO, "Libsumo." Documentation, accessed Nov. 11, 2026.
\bibitem{r6} Eclipse SUMO, "TripInfo." Metric definitions, accessed Nov. 11, 2026.
\end{thebibliography}
\end{document}